\input{../tex-inputs/latex-korrekturansicht-vorspann}

\section[Arthur Schnitzler an Berta Zuckerkandl, 20. 6. 1911]{L03950 Arthur Schnitzler an Berta Zuckerkandl, 20. 6. 1911}
\nopagebreak\mylabel{L03950v}
\rehead{ }\normalsize\beginnumbering\briefempfaengerindex{Zuckerkandl, Berta@\textsc{Zuckerkandl, Berta}!zzzSchnitzler, Arthur@\emph{von Arthur Schnitzler}!1911-06-201@{20. 6. 1911}|(be}
\toendnotes[C]{\smallbreak\pagebreak[2]}
\correspDesc{Versand  durch Arthur Schnitzler am 20. 6. 1911 in Wien
\newline{}Erhalt  durch Berta Zuckerkandl am 22. 6. 1911 in Paris}\toendnotes[C]{\smallbreak}
\Standort{DLA, HS.1985.1.2282.}
\physDesc{Brief, Durchschlag, 1 Blatt, 2 Seiten, 937 Zeichen
\newline{}Schreibmaschine
\newline{}Handschrift Schreibkraft: Bleistift, deutsche Kurrent (\noindent{}Monats- und Jahresangabe beim Datum: »Juni 1911.«, Korrekturen)
\newline{}Handschrift Arthur Schnitzler: roter Buntstift, lateinische Kurrent (\noindent{}Vermerk: »\uline{Zuckerkandl}« und »\textcolor{pink}{Fr{[}an{]}k{[}reich{]}}\oindex{Frankreich@\textbf{Frankreich}|pw}«, vier Unterstreichungen)}\toendnotes[C]{\smallbreak}
\pstart
           \raggedleft{}{\pb}20.\pend
           
\pstart\center{}Verehrte gnädige Frau.\pend\vspace{0.5em}
\pstart
           Hier sende ich Ihnen das \label{K_L03950-1v}\edtext{\textcolor{green}{Scenarium}\pwindex{Schnitzler, Arthur 15. 5. 1862 Wien – 21. 10. 1931 ebd.@\textsc{Schnitzler, Arthur} (15. 5. 1862 Wien – 21. 10. 1931 ebd.), \emph{Schriftsteller, Mediziner}!junge Medardus. (Zur Übersetzung ins Französische)@\strich\emph{Der junge Medardus. (Zur Übersetzung ins Französische)}|pwv}{}\ledrightnote{{$\rightarrow$}\emph{\textcolor{green}{Der junge Medardus. (Zur Übersetzung ins Französische)}}} des »\textcolor{green}{Medardus}\pwindex{Schnitzler, Arthur 15. 5. 1862 Wien – 21. 10. 1931 ebd.@\textsc{Schnitzler, Arthur} (15. 5. 1862 Wien – 21. 10. 1931 ebd.), \emph{Schriftsteller, Mediziner}!junge Medardus. Dramatische Historie in einem Vorspiel und fünf Aufzügen@\strich\emph{Der junge Medardus. Dramatische Historie in einem Vorspiel und fünf Aufzügen}|pw}{}\ledrightnote{\textcolor{green}{Der junge Medardus. Dramatische Historie in einem Vorspiel und fünf Aufzügen}}«}{\lemma{\textnormal{\emph{Scenarium des »Medardus«}}}\Cendnote{\textnormal{Die Beilage des Briefes ist nicht mit diesem überliefert. \textcolor{blue}{Schnitzler} vermerkte für den 16. 6. 1911 im \emph{\textcolor{green}{Tagebuch}\pwindex{Schnitzler, Arthur 15. 5. 1862 Wien – 21. 10. 1931 ebd.@\textsc{Schnitzler, Arthur} (15. 5. 1862 Wien – 21. 10. 1931 ebd.), \emph{Schriftsteller, Mediziner}!Tagebuch@\strich\emph{Tagebuch}|pwk}} die Arbeit am \textcolor{green}{Szenarium}\pwindex{Schnitzler, Arthur 15. 5. 1862 Wien – 21. 10. 1931 ebd.@\textsc{Schnitzler, Arthur} (15. 5. 1862 Wien – 21. 10. 1931 ebd.), \emph{Schriftsteller, Mediziner}!junge Medardus. (Zur Übersetzung ins Französische)@\strich\emph{Der junge Medardus. (Zur Übersetzung ins Französische)}|pwkv} und für den 20. 6. 1911 dessen
                  Absenden an \textcolor{blue}{Berta Zuckerkandl}\pwindex{Zuckerkandl, Berta 13.\,4.\,1864 Wien – 16.\,10.\,1945 Paris@\textsc{Zuckerkandl, Berta} (13.\,4.\,1864 Wien – 16.\,10.\,1945 Paris), \emph{Schriftstellerin, Journalistin, Übersetzerin}|pwk}. Ein Exemplar
                  ist im Nachlas \textcolor{blue}{Schnitzlers} in der \emph{Cambridge University Library} (A 94) erhalten. Es
                  ist getippt und enthält zusätzlich zum Titelblatt 35 Seiten. Siehe Arthur Schnitzler: \emph{Mikrofilme}, \url{https://schnitzler\_mikrofilme.acdh.oeaw.ac.at/1428195\_0004}.}}}\label{K_L03950-1} ein
               und hoffe, dass es Ihren Wünschen leidlich entsprechen dürfte. Die am Schlusse
               angefügte \label{K_L03950-2v}\edtext{Bemerkung}{\lemma{\textnormal{\emph{Bemerkung}}}\Cendnote{\textnormal{Es handelt sich um einen handschriftlich mit »\uline{Zusatz}« überschriebenen
                  Teil, der auf S. 34 und 35 des \textcolor{green}{Szenariums}\pwindex{Schnitzler, Arthur 15. 5. 1862 Wien – 21. 10. 1931 ebd.@\textsc{Schnitzler, Arthur} (15. 5. 1862 Wien – 21. 10. 1931 ebd.), \emph{Schriftsteller, Mediziner}!junge Medardus. (Zur Übersetzung ins Französische)@\strich\emph{Der junge Medardus. (Zur Übersetzung ins Französische)}|pwkv} zu finden ist: »Obwohl ich mir bewusst bin durch die Wahl
                     des Namens \textcolor{green}{Valois}\pwindex{Schnitzler, Arthur 15. 5. 1862 Wien – 21. 10. 1931 ebd.@\textsc{Schnitzler, Arthur} (15. 5. 1862 Wien – 21. 10. 1931 ebd.), \emph{Schriftsteller, Mediziner}!junge Medardus. Dramatische Historie in einem Vorspiel und fünf Aufzügen@\strich\emph{Der junge Medardus. Dramatische Historie in einem Vorspiel und fünf Aufzügen}|pwv} für die Herzogsfamilie
                     im ›\textcolor{green}{Jungen Medardus}\pwindex{Schnitzler, Arthur 15. 5. 1862 Wien – 21. 10. 1931 ebd.@\textsc{Schnitzler, Arthur} (15. 5. 1862 Wien – 21. 10. 1931 ebd.), \emph{Schriftsteller, Mediziner}!junge Medardus. Dramatische Historie in einem Vorspiel und fünf Aufzügen@\strich\emph{Der junge Medardus. Dramatische Historie in einem Vorspiel und fünf Aufzügen}|pw}‹ gegen die historische
                     Wahrheit im höheren Sinn\introOben{},\introOben{} gegen den Geist der
                     Geschichte \introOben{}mich\introOben{} nicht vergangen \introOben{}zu\introOben{} und keineswegs die Grenzen überschritten \introOben{}zu\introOben{} habe\introOben{}n\introOben{}, die dem
                     Dichter einer dramatischen Historie bei der
                     Verwendung geschichtlicher Fakten und Namen
                     gesetzt werden müssen, so stelle ich es doch,
                     um wohlfeilen Erinnerungen der Kritik zuvorzukommen dem französischen Uebersetzer 
                     meines \textcolor{green}{Stückes}\pwindex{Schnitzler, Arthur 15. 5. 1862 Wien – 21. 10. 1931 ebd.@\textsc{Schnitzler, Arthur} (15. 5. 1862 Wien – 21. 10. 1931 ebd.), \emph{Schriftsteller, Mediziner}!junge Medardus. Dramatische Historie in einem Vorspiel und fünf Aufzügen@\strich\emph{Der junge Medardus. Dramatische Historie in einem Vorspiel und fünf Aufzügen}|pwv} anheim den Namen \textcolor{green}{Valois}\pwindex{Schnitzler, Arthur 15. 5. 1862 Wien – 21. 10. 1931 ebd.@\textsc{Schnitzler, Arthur} (15. 5. 1862 Wien – 21. 10. 1931 ebd.), \emph{Schriftsteller, Mediziner}!junge Medardus. Dramatische Historie in einem Vorspiel und fünf Aufzügen@\strich\emph{Der junge Medardus. Dramatische Historie in einem Vorspiel und fünf Aufzügen}|pwv} durch
                     irgend einen andern, meinetwegen selbst erfundenen Aristokratennamen zu ersetzen. Wird
                     eine entferntere \introOben{}(\introOben{}vielleicht gar illegitime\introOben{})\introOben{}
                     Verwandtschaft der in \textcolor{pink}{Wien}\oindex{Wien@\textbf{Wien}, \emph{Verwaltungsgebiet}|pw} lebenden Emigrantenfamilie zu dem bourbinischen Königshaus
                     supponiert als in meinem \textcolor{green}{Text}\pwindex{Schnitzler, Arthur 15. 5. 1862 Wien – 21. 10. 1931 ebd.@\textsc{Schnitzler, Arthur} (15. 5. 1862 Wien – 21. 10. 1931 ebd.), \emph{Schriftsteller, Mediziner}!junge Medardus. Dramatische Historie in einem Vorspiel und fünf Aufzügen@\strich\emph{Der junge Medardus. Dramatische Historie in einem Vorspiel und fünf Aufzügen}|pwv}, so tritt das
                     Närrische im Wesen des Herzogs, das phantastische in dem der Prinzessin am Ende gar zum Vorteil der theatralischen Wirkung noch
                     stärker in der Erscheinung als in der jetzt
                     vorliegenden Fassung.« Arthur Schnitzler: \emph{Mikrofilme}, \url{https://schnitzler\_mikrofilme.acdh.oeaw.ac.at/1428195\_0038}.}}}\label{K_L03950-2} dürfte wohl der Erwägung wert sein. Hoffentlich habe ich bald das
               Vergnügen mehr von der Sache zu hören und danke Ihnen für Ihr Interesse sowohl als
               für Ihre Bemühungen.\pend
           
\pstart
           Der Ordnung wegen möchte ich es hier auch noch schriftlich niederlegen, dass ich für
               den Fall, dass wir zu einem Resultate kommen sollten, auf die Hälfte der Tantiemen
               und der eventuellen anderen Honorare, Buchausgabe und dergl., Anspruch mache. Wegen
               der Verteilung der andern Hälfte werden sie sich wohl selbst, verehrte gnädige Frau,
               mit Ihrer präsumptiven Hilfskraft ins Einvernehmen setzen. Ich schreibe an den \textcolor{brown}{Verlag}\orgindex{S. Fischer Verlag@S. Fischer Verlag|pwv}{}\ledrightnote{{$\rightarrow$}\emph{\textcolor{brown}{S. Fischer Verlag}}} wegen der an Sie zu
               sendenden \textcolor{green}{Exemplare}\pwindex{Schnitzler, Arthur 15. 5. 1862 Wien – 21. 10. 1931 ebd.@\textsc{Schnitzler, Arthur} (15. 5. 1862 Wien – 21. 10. 1931 ebd.), \emph{Schriftsteller, Mediziner}!junge Medardus. Dramatische Historie in einem Vorspiel und fünf Aufzügen@\strich\emph{Der junge Medardus. Dramatische Historie in einem Vorspiel und fünf Aufzügen}|pwv}{}\ledrightnote{{$\rightarrow$}\emph{\textcolor{green}{Der junge Medardus. Dramatische Historie in einem Vorspiel und fünf Aufzügen}}}. Es stehen
               Ihnen natürlich nach Bedarf auch weitere zur Verfügung.\pend
           
\pstart
           {\pb}Mit herzlichen Grüssen{\\[\baselineskip]} Ihr sehr ergebener\pend
           \leftskip=0em{}{\vspace{1\baselineskip}}
\pstart
           \noindent{}Frau Berta Zuckerkandl, \label{K_L03950-3v}\edtext{\textcolor{pink}{Wien}\oindex{Wien@\textbf{Wien}, \emph{Verwaltungsgebiet}|pw}{}\ledrightnote{\textcolor{pink}{Wien}}}{\lemma{\textnormal{\emph{Wien}}}\Cendnote{\textnormal{Aus dem \emph{\textcolor{green}{Tagebuch}\pwindex{Schnitzler, Arthur 15. 5. 1862 Wien – 21. 10. 1931 ebd.@\textsc{Schnitzler, Arthur} (15. 5. 1862 Wien – 21. 10. 1931 ebd.), \emph{Schriftsteller, Mediziner}!Tagebuch@\strich\emph{Tagebuch}|pwk}}eintrag zum 20. 6. 1911 geht hervor, dass \textcolor{blue}{Schnitzler} den vorliegenden Brief nicht
                     nach \textcolor{pink}{Wien}\oindex{Wien@\textbf{Wien}, \emph{Verwaltungsgebiet}|pwk} sondern nach \textcolor{pink}{Paris}\oindex{Paris@\textbf{Paris}, \emph{Hauptstadt}|pwk} sandte, wohin \textcolor{blue}{Zuckerkandl}\pwindex{Zuckerkandl, Berta 13.\,4.\,1864 Wien – 16.\,10.\,1945 Paris@\textsc{Zuckerkandl, Berta} (13.\,4.\,1864 Wien – 16.\,10.\,1945 Paris), \emph{Schriftstellerin, Journalistin, Übersetzerin}|pwk} am 17. 6. 1911 abgereist sein dürfte, vgl. Berta Zuckerkandl an Arthur Schnitzler, [13. 6. 1911?].}}}\label{K_L03950-3}.\pend
           \selectlanguage{ngerman}\endnumbering\briefempfaengerindex{Zuckerkandl, Berta@\textsc{Zuckerkandl, Berta}!zzzSchnitzler, Arthur@\emph{von Arthur Schnitzler}!1911-06-201@{20. 6. 1911}|)be}\mylabel{L03950h}
\begin{anhang}
\end{anhang}\newcommand{\dateiname}{L03950}\newcommand{\titel}{Arthur Schnitzler an Berta Zuckerkandl, 20. 6. 1911}\newcommand{\editorInnen}{Herausgegeben von Jahnke, SelmaMüller, Martin Anton}\input{../tex-inputs/latex-korrekturansicht-abspann}
